\clearpage
\fancyhead[L]{\small PASQUICH ASTRONOMIA SPHAERICA.}
\SourceLeaf{540}
\begin{center}\small Göttingische gelehrte Anzeigen. Stück 144. Seite 1433..1436. 1811 September 9.\\[3pt]\rule{.72\linewidth}{.3pt}\end{center}

{\rightskip=0pt plus 1em \emph{Epitome elementorum astronomiae sphaerico-calculatoriae auctore} \textsc{Joanne Pasquich}, directore observatorii astronomici regiae universitatis hungaricae. \emph{Pars prima. Elementa {\addfontfeatures{LetterSpace=10}theoretica} astronomiae sphaerico-calculatoriae. Pars secunda, Elementa {\addfontfeatures{LetterSpace=10}practica} astronomiae sphaerico-calculatoriae.} 1811. 160 und 166 Seiten in Quart. Wien, Schaumann u. Comp.\par}

Der Verf. erklärt in der Vorrede, dass diese Schrift bei Gelegenheit von astronomischen, auf der Ofener Universität gehaltenen, Vorlesungen entstanden, dass sie nicht für geübte Astronomen, sondern nur für erste Anfänger bestimmt sei, die dadurch nicht sowohl eine vollständige Anleitung zu allen mannigfaltigen astronomischen Vorrichtungen erhalten, als vielmehr nur zu dem dort bisher zu sehr vernachlässigten Studium dieser Wissenschaft angelockt und aufgemuntert werden sollen, um demnächst einen weiter gehenden Unterricht zu empfangen, und dass man hiermit die sonst auffallende Dürftigkeit des Buches zu entschuldigen habe. Es gereicht dem Verf. zur Ehre, dass er diesen Zweck nicht durch ein oberflächliches Abschöpfen unterhaltender Resultate zu erreichen sucht (wodurch das ernste Studium nicht gefördert wird), sondern von einer gründlichen Vorbereitung in den Anfangsgründen ausgeht. Der erste, theoretische, Theil enthält in fünf Abschnitten folgende Lehren. Der \emph{erste} Abschnitt unter dem Titel: Grundbegriff und Grundgesetze der täglichen Bewegung der Gestirne und der jährlichen Bewegung der Sonne, zerfällt in vier Kapitel, welche die ersten Grundbegriffe von der Sphäre und den darauf gezogenen Kreisen, so wie die trigonometrischen Formeln zur Bestimmung der relativen Lage der Gestirne gegen jene Kreise, darlegen. Der gleichfalls in vier Kapitel zerfallende \emph{zweite} Abschnitt beschäftigt sich mit den verschiedenen Arten von Zeit und dem darauf sich beziehenden Gebrauch der geraden Aufsteigungen, so wie mit den Unterschieden der Meridiane. Bei der sonst überall herrschenden Gründlichkeit des Vortrags könnte man doch wünschen, dass hier bei einem oder andern Punkte noch etwas tiefer eingedrungen wäre. So ist z. B. der Begriff von mittlerer Sonnenzeit als \emph{Zeitdauer} zwar ganz klar, aber der Begriff mittlere Sonnenzeit als \emph{Zeitpunkt} bleibt noch schwankend, da man nicht sieht, von wo aus man die mittlere Bewegung der Sonne mit der wahren zugleich auslaufend annehmen soll. Der \emph{dritte} Abschnitt handelt in drei Kapiteln von den Beziehungen zwischen der Bewegung der Sonne in der Ekliptik und ihren Rectascensionen, Declinationen und Positionswinkeln; von den Höhen, Stundenwinkeln und Azimuthen der Gestirne. Im \emph{vierten} Abschnitte wird in vier Kapiteln von der Gestalt und Grösse der Erde, von der Parallaxe und von der astronomischen Strahlenbrechung gehandelt. Im \emph{fünften} Abschnitte lehrt der Verf., in fünf Kapiteln, die ersten Gründe der elliptischen Bewegung der Planeten und der davon abhängenden Erscheinungen, die Wirkungen der Präcession, Nutation und Aberration, womit der theoretische Theil beendigt wird. Man sieht aus dieser Inhaltsanzeige, dass die Absicht des Verf. mehr dahin gegangen ist, ein Repertorium brauchbarer Rechnungsformeln zu geben, als eine Ordnung des Vortrags zu wählen, wodurch man zu einer belehrenden und befriedigenden Einsicht in die Art, wie die Wahrheiten der Astronomie gefunden sind oder gefunden werden könnten, ohne mündliche Nachhülfe gelangen könnte, welcher also Vieles überlassen bleiben muss.

Der \emph{zweite Theil} des Werks ist keinesweges als ein vollständiges Handbuch der rechnenden Astronomie anzusehen, sondern beschränkt sich nur auf solche practische Vorschriften, deren etwa der beobachtende Astronom am häufigsten bedarf. Er enthält drei Abschnitte. Im \emph{ersten} wird die Einrichtung astronomischer Tafeln im Allgemeinen, und der Sonnentafeln besonders, der Gebrauch der astronomi\SourceLeaf{541}schen Ephemeriden, der Fixsternverzeichnisse und die Berechnung der Nutation und Aberration practisch erläutert. Im \emph{zweiten} Abschnitte wird die Vergleichung der verschiedenen astronomischen Zeiten unter einander und die Zeitbestimmungen durch beobachtete Culminationen und correspondirende Höhen gelehrt. Im \emph{dritten} Abschnitt werden vornehmlich der Gebrauch beobachteter Circummeridianhöhen zu Breitenbestimmungen gezeigt, und die Art, wie die Stellungen der Himmelskörper aus den Beobachtungen abzuleiten sind, beschrieben.

Zu obigem Werke gehört noch: Appendix ad Joannis Pasquich epitomen elementorum practicorum astronomiae sphaerico-calculatoriae, complectens tabulas auxiliares. Viennae apud Schaumburg et Soc. 1811. 42 Seiten in Quart. Diese Tafeln sind: Die \textsc{Delambre}'schen Sonnentafeln, abgekürzt; die allgemeine Tafel für die Mittagsverbesserung nach der gewöhnlichen Einrichtung; eine Hülfstafel zur Reduction ausser dem Meridian gemessener Zenithdistanzen; eine andere Hülfstafel zur Reduction der Zenithdistanzen des Polarsterns; eine Tafel zur Bestimmung der Präcession in gerader Aufsteigung und Abweichung; Tafeln für die Refraction, Aberration und Nutation; der neueste \textsc{Piazzi}'sche Catalog von 121 Fundamentalsternen. Rec. hätte gewünscht, dass einigen dieser Tafeln, deren Einrichtung für sich nicht verständlich ist (wie den Tafeln für die Reduction der Zenithdistanzen), die Erklärung beigefügt wäre, da vielleicht Mancher sich die Tafeln ohne die Epitome anschaffen möchte.
